\newpage
\section{Discussion and Limitations}\label{ch:discussion}
\subsection{What the Models Learn}
All learned models exceed the mean Macro-F1 of the uniform-random baseline, but this does not imply that every model is reliably better in a statistical sense. Label distributions vary across regimes: the Neutral share is 64.2\% in Low-vol, 54.9\% in Mid-vol and 28.8\% in High-vol across the full set of rows. This is an association between recent market conditions and the next-day class distribution; it does not establish that most predictive skill comes from forecasting volatility magnitude rather than direction. The price-only model is already competitive, and adding a regime label to the price-plus-embedding model does not produce a reliable overall improvement.

\subsection{Why the Embedding Hurts}
The 768-dimensional embedding must be mapped to a useful representation from about 1{,}460 training rows. The train--test gaps in Table~\ref{tab:overfit} (0.138 for M2, 0.319 for M4, against 0.011 for M1) are the signature of overfitting, and the significant decrease when the embedding is added to the price model is consistent with it. Two further factors may contribute. First, the daily mean-pooled embedding varies little from one day to the next (cosine similarity 0.963 between consecutive days and 0.951 between random days), so the useful contrast is small relative to the 768 dimensions. Second, the headline supply and tone change across the decade (49 headlines per day in training versus 26 in test after capping; mean sentiment $-0.068$ versus $-0.184$), so a high-capacity news model can learn era-specific patterns that do not transfer. This report did not test these explanations separately; they are hypotheses consistent with the data.

\subsection{Why Protocol Effects Are Small}
The regime-only model scores higher under both looser protocols, consistent with a slowly varying regime assignment being sensitive to random splitting and to fitting clusters on the full sample. However, P1 uses a different random test set for each seed, while P2/P3 use the chronological test set. Consequently, the protocol gaps combine protocol changes with test-sample composition and run variability; they do not isolate a causal effect of leakage. These results apply only to the specific split and preprocessing changes examined here. They do not rule out other forms of leakage, including use of information unavailable at prediction time.

\subsection{The Stress Window}
News did not add reliable value during the February--April 2020 episode in this data, although headline sentiment clearly moved. With 50 test rows the analysis has low power, so the result is best read as ``no large effect detected'' and not as ``no effect''. Because this study uses one benchmark, a small stress subset and date-only headlines, it cannot settle whether news is generally more valuable during volatile markets. Differences in assets, label construction, data timing and evaluation protocols prevent direct comparison with published point estimates.

\subsection{Limitations}
\begin{itemize}
\item One asset (SPY), one benchmark and one decade; results may not transfer to other markets.
\item Only 317 test rows (125, 113 and 79 by regime; 50 in the stress window): intervals are wide and small effects cannot be detected. Test rows are overlapping, autocorrelated daily observations, so bootstrap intervals that treat them as independent are probably somewhat too narrow.
\item Headlines carry a date but no publication time. A headline dated $t$ precedes the predicted move only if it was published by the close of day $t$; this could not be verified.
\item Headline volume and tone differ across the three periods, so train and test distributions differ.
\item Regimes are data-driven clusters (three clusters on rolling 20-day return and log-volatility), not validated economic states; the Mid-vol and High-vol regimes switch frequently (median run length 4 days).
\item The three protocols were evaluated on different test samples (P1 uses a random test set), so protocol gaps combine leakage and sampling differences.
\item About 24 comparisons were examined; isolated significant intervals are hypotheses. M6 and the random-guess baseline were added after seeing pilot results (Appendix~\ref{app:transparency}).
\item Fifteen expected trading days are missing from the recovered daily series and were not interpolated.
\item No transaction costs, trading rule or profitability analysis is attempted, and no causal claim is made.
\end{itemize}
